
\section{非光滑丛}\label{sec:C-6}

在若干应用中，我们必须考虑光滑性较弱的丛。例如，第 3 章中需要处理 $W^{\ell,p}$ 类的丛以及 $W^{\ell-1,p}$ 类的联络，其中 $\ell$ 是整数，$p>1$，且 $\ell p>2$。更明确地，令 $(E,F)$ 是定义~\ref{def:C-3-4} 所述的、带边紧黎曼曲面 $(\Sigma,j)$ 上的\BookRecordIndex{T0080}丛对，并设 $\langle\cdot,\cdot\rangle$ 是 $E$ 上的埃尔米特内积（即满足 $J^*=-J$ 的内积），且 $JF=F^\perp$。我们称三元组
\[
  (E,F,\langle\cdot,\cdot\rangle)
\]
为\emph{\BookRecordIndex{T0051}埃尔米特丛对}。\BookRecordIndex{T0051}埃尔米特丛对 $(E,F,\langle\cdot,\cdot\rangle)$ 具有一组酉局部平凡化
\[
  \phi_\alpha:\pi^{-1}(U_\alpha)\longrightarrow U_\alpha\times\mathbb C^n,
\]
其中 $\{U_\alpha\}_\alpha$ 是 $\Sigma$ 的适当开覆盖，并且
\[
  \phi_\alpha\bigl(\pi^{-1}(U_\alpha)\cap F\bigr)
  =(U_\alpha\cap\partial\Sigma)\times\mathbb R^n.
\]
转移函数
\[
  g_{\beta\alpha}:U_\alpha\cap U_\beta\longrightarrow U(n)
\]
由
\[
  (z,g_{\beta\alpha}(z)v):=\phi_\beta\circ\phi_\alpha^{-1}(z,v)
\]
定义，且满足
\begin{equation}\label{eq:C-6-1}
  g_{\gamma\beta}g_{\beta\alpha}=g_{\gamma\alpha},\qquad
  g_{\alpha\alpha}=\mathrm{id},\qquad
  g_{\beta\alpha}(U_\alpha\cap U_\beta\cap\partial\Sigma)\subset O(n).
\end{equation}
反过来，每一组满足~\eqref{eq:C-6-1} 的转移函数 $g_{\beta\alpha}:U_\alpha\cap U_\beta\to U(n)$ 都决定一个\BookRecordIndex{T0051}埃尔米特丛对 $(E,F,\langle\cdot,\cdot\rangle)$。（试证明！）在这些局部平凡化下，\BookRecordIndex{T0052}埃尔米特联络 $\nabla$ 具有形式
\[
  (\nabla\xi)_\alpha=d\xi_\alpha+A_\alpha\xi_\alpha
  \in\Omega^1(U_\alpha,\mathbb C^n),
\]
其中 $\xi_\alpha:U_\alpha\to\mathbb C^n$ 和 $(\nabla\xi)_\alpha\in\Omega^1(U_\alpha,\mathbb C^n)$ 由
\[
  \phi_\alpha\circ\xi(z)=:(z,\xi_\alpha(z)),\qquad
  \phi_\alpha^*(\nabla\xi)_\alpha:=\nabla\xi|_{U_\alpha}
\]
定义。$1$-形式 $A_\alpha\in\Omega^1(U_\alpha,\mathfrak u(n))$ 称为\emph{\BookRecordIndex{T0223}联络势}。注意 $\xi_\alpha$ 和 $A_\alpha$ 满足
\[
  g_{\beta\alpha}\xi_\alpha=\xi_\beta,\qquad
  A_\alpha=g_{\beta\alpha}^{-1}dg_{\beta\alpha}
    +g_{\beta\alpha}^{-1}A_\beta g_{\beta\alpha}.
\]

\begin{definition}\label{def:C-6-1}
令 $\ell$ 为整数，$p>1$，且 $\ell p>2$。若一个（连续的）\BookRecordIndex{T0051}埃尔米特丛对 $(E,F)$ 具有上述平凡化，且其转移函数 $g_{\alpha\beta}$ 属于 $W^{\ell,p}$ 类，则称它\emph{属于 $W^{\ell,p}$ 类}。若联络 $\nabla$ 的\BookRecordIndex{T0223}联络势 $A_\alpha$ 属于 $W^{\ell-1,p}$ 类，则称该联络属于 $W^{\ell-1,p}$ 类。$W^{\ell,p}$ 类\BookRecordIndex{T0051}埃尔米特丛对上的\emph{$W^{\ell-1,p}$ 类\BookRecordIndex{T0210}柯西--黎曼算子}是形如 $D=\bar\partial^\nabla$ 的算子，其中 $\nabla$ 是 $W^{\ell-1,p}$ 类\BookRecordIndex{T0052}埃尔米特联络。
\end{definition}

下面的习题表明，每个 $W^{\ell,p}$ 类\BookRecordIndex{T0051}埃尔米特丛对 $(E,F)$ 都与某个光滑\BookRecordIndex{T0051}埃尔米特丛对 $(E_0,F_0)$ 在 $W^{\ell,p}$ 类中等距同构。

\begin{exercise}[{$W^{\ell,p}$ 类丛}]\label{ex:C-6-2}
令 $(E,F)$ 是 $W^{\ell,p}$ 类\BookRecordIndex{T0051}埃尔米特丛对，局部平凡化为 $\phi_\alpha:\pi^{-1}(U_\alpha)\to U_\alpha\times\mathbb C^n$。
\begin{enumerate}[label=(\roman*)]
\item 选择一族光滑\BookRecordIndex{T0188}截断函数 $\sigma_\alpha:\Sigma\to[0,1]$，使 $\operatorname{supp}\sigma_\alpha\subset U_\alpha$，且 $\rho_\alpha:=\sigma_\alpha^2$ 构成单位分解。定义嵌入
\[
  E\longrightarrow\Sigma\times\mathbb C^N:\quad
  (z,e)\longmapsto(z,\Phi(z)e),\qquad
  \Phi(z)e:=\{\sigma_\alpha(z)\Phi_\alpha(z)e\}_\alpha,
\]
其中 $\Phi_\alpha(z):E_z\to\mathbb C^n$ 由 $\phi_\alpha(z,e)=(z,\Phi_\alpha(z)e)$ 给出。证明 $\Phi(z)$ 是线性等距嵌入，并且对 $z\in\partial\Sigma$ 有 $\Phi(z)F_z=\operatorname{im}\Phi(z)\cap\mathbb R^N$。
\item 令 $\Phi$ 如 (i) 所述，并定义 $\Pi:\Sigma\to\mathbb C^{N\times N}$ 为
\[
  \Pi(z):=\Phi(z)\bigl(\Phi(z)^*\Phi(z)\bigr)^{-1}\Phi(z)^*.
\]
证明 $\Pi(z)$ 是 $\mathbb C^N$ 到 $\Phi(z)$ 的像上的正交投影。证明 $\Pi\in W^{\ell,p}(\Sigma,\mathbb C^{N\times N})$，且 $\Pi(\partial\Sigma)\subset\mathbb R^{N\times N}$。
\item 令 $\Pi_0:\Sigma\to\mathbb C^{N\times N}$ 是光滑函数，且满足
\[
  \Pi_0(z)^2=\Pi_0(z)=\Pi_0(z)^*,\qquad
  \Pi_0(\partial\Sigma)\subset\mathbb R^{N\times N},\qquad
  \max_{z\in\Sigma}|\Pi(z)-\Pi_0(z)|<1.
\]
证明对每个 $z\in\Sigma$，都有 $\operatorname{rank}\Pi_0(z)=\operatorname{rank}\Pi(z)=n$。
\item 对 $z\in\Sigma$ 令 $E_{0z}:=\operatorname{im}\Pi_0(z)$，对 $z\in\partial\Sigma$ 令 $F_{0z}:=\operatorname{im}\Pi_0(z)\cap\mathbb R^N$，由此定义 $\Sigma$ 上的光滑\BookRecordIndex{T0051}埃尔米特丛对 $(E_0,F_0)$。证明映射 $\Psi_0:(E_0,F_0)\to(E,F)$，$\Psi_0:=\Phi^{-1}\Pi$，是 $W^{\ell,p}$ 类丛同构。证明
\[
  \Phi_0:=\Psi_0(\Psi_0^*\Psi_0)^{-1/2}
\]
是 $W^{\ell,p}$ 类埃尔米特丛同构。
\end{enumerate}
\end{exercise}

\begin{corollary}\label{cor:C-6-3}
令 $\ell$ 为正整数，$p>1$，且 $\ell p>2$。设 $(E,F)$ 是 $W^{\ell,p}$ 类\BookRecordIndex{T0051}埃尔米特丛对，$D$ 是 $E$ 上 $W^{\ell-1,p}$ 类\BookRecordIndex{T0210}柯西--黎曼算子。那么定理~\ref{thm:C-1-10} 的断言对满足 $k\in\{1,\ldots,\ell\}$、$q>1$ 且 $k-2/q\le\ell-2/p$ 的 $k,q$ 仍成立。
\end{corollary}
\begin{proof}
由习题~\ref{ex:C-6-2}，可以假设\BookRecordIndex{T0051}埃尔米特丛对 $(E,F)$ 是光滑的。于是结论由定理~\ref{thm:C-1-10} 得到。
\end{proof}
